\chapter{INTRODUCTION}

\section{Introduction}

\subsection{Problem Statement}
Sign language is the primary means of communication for approximately 466 million people worldwide who experience significant hearing loss \cite{who2021deafness}. Unlike spoken language, sign language encodes meaning simultaneously through hand shape, movement trajectory, body posture, and -- for certain vocabulary items -- facial expression \cite{ong2005automatic}. This multi-channel, spatiotemporal complexity makes automated Sign Language Recognition (SLR) one of the most challenging tasks in computer vision and human-computer interaction.

Despite rapid advances in deep learning for gesture and pose estimation \cite{koller2020weakly}, Vietnamese Sign Language (VSL) research remains far behind resource-rich languages such as American Sign Language (ASL) or German Sign Language (DGS). The root cause is not algorithmic: state-of-the-art models trained on ASL datasets achieve above 90\% word-level accuracy \cite{li2020wlasl}. The bottleneck is \textbf{data}. Training a competitive isolated SLR model requires tens of thousands of annotated video clips \cite{koller2020weakly}, yet no publicly accessible, community-driven collection infrastructure exists for VSL.

A preliminary study conducted at Can Tho University, involving three active VSL research groups and five deaf community volunteers, revealed three concrete operational pain-points:
\begin{enumerate}
    \item \textbf{Data loss under network failure.} Researchers uploading large video batches directly to Google Drive from the browser lose all progress when the connection drops mid-upload, requiring complete re-upload with no partial-credit recovery.
    \item \textbf{Undetected cross-group duplication.} Multiple research groups independently record the same vocabulary items, wasting collection effort and making dataset merger impossible without manual deduplication.
    \item \textbf{Absence of a canonical taxonomy.} Without a shared vocabulary library, each annotator invents ad-hoc label names, making cross-group dataset merging nearly impossible and hindering reproducibility.
\end{enumerate}

\subsection{Proposed Solution}
To address the aforementioned challenges, this research proposes \textbf{CTU-SignBridge} -- a purpose-built, self-hosted Sign Language data collection and management platform. CTU-SignBridge aims to provide a centralized infrastructure for capturing, managing, and standardizing VSL datasets. By incorporating real-time landmark extraction on the client-side, enforcing a canonical taxonomy through a structured database, and utilizing asynchronous task queues for resilient data synchronization, the system provides a robust solution for large-scale dataset creation.

\subsection{Research Objectives}
The primary goal of this research is to design and implement \textbf{CTU-SignBridge}, a comprehensive web-based platform tailored for the scalable collection, management, and processing of Vietnamese Sign Language (VSL) datasets. To achieve this overarching goal, the thesis pursues the following specific objectives:

\begin{enumerate}
    \item \textbf{Develop a Multi-Tenant Web Platform Architecture:} Design a robust, multi-tenant Workspace--Project architecture governed by strict Role-Based Access Control (RBAC) powered by Casbin and authentication managed by Keycloak. This ensures secure, isolated, and collaborative data management environments.
    \item \textbf{Optimize Storage through In-Browser Landmark Extraction:} Implement a real-time, client-side data processing pipeline utilizing MediaPipe Holistic. By extracting and storing essential skeletal keypoints (landmarks) as highly compressed \texttt{.npz} arrays rather than raw RGB video files, the system aims to achieve over a 90\% reduction in storage footprint and network bandwidth.
    \item \textbf{Design a Relational VSL Taxonomy:} Construct a structured, multi-dimensional database schema (Star Schema) capable of capturing the complex dialectal variations of VSL.
    \item \textbf{Ensure High Reliability and Data Durability:} Architect an asynchronous backend processing pipeline leveraging Celery and Redis message brokers, combined with MinIO for scalable object storage, guaranteeing zero data loss during cloud synchronization.
    \item \textbf{Evaluate System Performance and Scalability:} Conduct rigorous empirical validation to ensure the platform meets production-grade Key Performance Indicators (KPIs) via load testing and isolation validation.
\end{enumerate}

\subsection{Scope and Limitations}
This project covers the \textbf{data collection and management layer} of a full VSL recognition pipeline. The following items are explicitly \textit{out of scope}:
\begin{itemize}
    \item Training or evaluating sign language recognition models (only the dataset infrastructure is built).
    \item Automatic real-time sign recognition (no inference engine in this release).
    \item Facial expression capture: the SIGN\_FEATURES table reserves the \texttt{requires\_face\_expression} column, but the corresponding capture pipeline is disabled in the current release.
    \item Deployment at internet scale (benchmarks are conducted on a 4-core, 16\,GB RAM test node using Docker Compose; Kubernetes-based horizontal scaling is future work).
\end{itemize}

\subsection{Practical Significance}
The outcome of this research provides a vital infrastructure foundation for the artificial intelligence and deaf communities in Vietnam. By drastically reducing the technical friction involved in recording and managing gesture data, CTU-SignBridge accelerates the creation of the first large-scale VSL dataset. In the long term, this data will enable the training of robust VSL translation models, thereby fostering accessibility and bridging the communication gap for the deaf community.

\section{Related Works}

\subsection{Domestic Research}
Within Vietnam, research on VSL recognition has largely focused on algorithm development rather than data infrastructure. Studies from local universities typically rely on small-scale, privately collected datasets containing limited vocabularies (e.g., 50 to 100 words), recorded in controlled environments without a unified taxonomy. A robust, open-source platform to aggregate these fragmented efforts has not yet been realized.

\subsection{Foreign Research}
Globally, platforms for sign language datasets are more mature. Major datasets such as WLASL (Word-Level American Sign Language) and AUTSL (Ankara University Turkish Sign Language Dataset) provide extensive public data but lack the dynamic collection interfaces for continuous community contribution. Generic data annotation tools like Label Studio or CVAT are highly capable but lack the specialized features needed for sign language (e.g., built-in MediaPipe landmark extraction, multi-tenant academic hierarchies, and sign-specific taxonomy structures).

\section{Thesis Organization}

The remainder of this thesis is organized as follows:
\begin{itemize}
    \item \textbf{Chapter 2 -- Theoretical Background:} Reviews related work on SLR datasets, MediaPipe Holistic, Star Schema taxonomy design, multi-tenancy architecture, IAM (Keycloak and Casbin), and asynchronous task processing.
    \item \textbf{Chapter 3 -- System Analysis and Design:} Presents the Clean Architecture, the database ERD, IAM design, and the asynchronous processing pipeline.
    \item \textbf{Chapter 4 -- Implementation and Evaluation:} Reports implementation details, tech stack usage, benchmark results for API latency, storage efficiency, and multi-tenant isolation.
    \item \textbf{Chapter 5 -- Conclusions and Future Work:} Summarizes contributions, evaluates KPI achievement, and proposes directions for future work.
\end{itemize}
